\section{AMC 8 1990}
        \begin{enumerate}
        	\item (\textit{AMC 8 1990}) What is the smallest sum of two $3$-digit numbers that can be obtained by placing each of the six digits $4,5,6,7,8,9$ in one of the six boxes in this addition problem?


 
                    \begin{figure}[H]
                        \centering
                        \includegraphics[width=0.5\textwidth]{Problem/AMC_8_AJHSME_1990_1_Problem_Figure_1.png}
                    \end{figure}
                    
$\text{(A)}\ 947 \qquad \text{(B)}\ 1037 \qquad \text{(C)}\ 1047 \qquad \text{(D)}\ 1056 \qquad \text{(E)}\ 1245$


	\item (\textit{AMC 8 1990}) Which digit of $.12345$, when changed to $9$, gives the largest number?


 $\text{(A)}\ 1 \qquad \text{(B)}\ 2 \qquad \text{(C)}\ 3 \qquad \text{(D)}\ 4 \qquad \text{(E)}\ 5$


	\item (\textit{AMC 8 1990}) What fraction of the square is shaded?


 
                    \begin{figure}[H]
                        \centering
                        \includegraphics[width=0.5\textwidth]{Problem/AMC_8_AJHSME_1990_3_Problem_Figure_1.png}
                    \end{figure}
                    
$\text{(A)}\ \frac{1}{3} \qquad \text{(B)}\ \frac{2}{5} \qquad \text{(C)}\ \frac{5}{12} \qquad \text{(D)}\ \frac{3}{7} \qquad \text{(E)}\ \frac{1}{2}$


	\item (\textit{AMC 8 1990}) Which of the following could \textbf{not} be the unit's digit [one's digit] of the square of a whole number?


 $\text{(A)}\ 1 \qquad \text{(B)}\ 4 \qquad \text{(C)}\ 5 \qquad \text{(D)}\ 6 \qquad \text{(E)}\ 8$


	\item (\textit{AMC 8 1990}) Which of the following is closest to the product $(.48017)(.48017)(.48017)$?


 $\text{(A)}\ 0.011 \qquad \text{(B)}\ 0.110 \qquad \text{(C)}\ 1.10 \qquad \text{(D)}\ 11.0 \qquad \text{(E)}\ 110$


	\item (\textit{AMC 8 1990}) Which of these five numbers is the largest?


 $\text{(A)}\ 13579+\frac{1}{2468} \qquad \text{(B)}\ 13579-\frac{1}{2468} \qquad \text{(C)}\ 13579\times \frac{1}{2468}$
$\text{(D)}\ 13579\div \frac{1}{2468} \qquad \text{(E)}\ 13579.2468$


	\item (\textit{AMC 8 1990}) When three different numbers from the set $\{ -3, -2, -1, 4, 5 \}$ are multiplied, the largest possible product is 


 $\text{(A)}\ 10 \qquad \text{(B)}\ 20 \qquad \text{(C)}\ 30 \qquad \text{(D)}\ 40 \qquad \text{(E)}\ 60$


	\item (\textit{AMC 8 1990}) A dress originally priced at $80$ dollars was put on sale for $25\%$ off.  If $10\%$ tax was added to the sale price, then the total selling price (in dollars) of the dress was


 $\text{(A)}\ \text{45 dollars} \qquad \text{(B)}\ \text{52 dollars} \qquad \text{(C)}\ \text{54 dollars} \qquad \text{(D)}\ \text{66 dollars} \qquad \text{(E)}\ \text{68 dollars}$


	\item (\textit{AMC 8 1990}) The grading scale shown is used at Jones Junior High.  The fifteen scores in Mr. Freeman's class were: 
 \[\begin{tabular}[t]{lllllllll} 89, & 72, & 54, & 97, & 77, & 92, & 85, & 74, & 75, \\ 63, & 84, & 78, & 71, & 80, & 90. & & & \\ \end{tabular}\]


 In Mr. Freeman's class, what percent of the students received a grade of C?


 
 \[\boxed{\begin{tabular}[t]{cc} A: & 93 - 100 \\ B: & 85 - 92 \\ C: & 75 - 84 \\ D: & 70 - 74 \\ F: & 0 - 69  \end{tabular}}\]
$\text{(A)}\ 20\% \qquad \text{(B)}\ 25\% \qquad \text{(C)}\ 30\% \qquad \text{(D)}\ 33\frac{1}{3}\% \qquad \text{(E)}\ 40\%$


	\item (\textit{AMC 8 1990}) On this monthly calendar, the date behind one of the letters is added to the date behind $\text{C}$.  If this sum equals the sum of the dates behind $\text{A}$ and $\text{B}$, then the letter is


 
                    \begin{figure}[H]
                        \centering
                        \includegraphics[width=0.5\textwidth]{Problem/AMC_8_AJHSME_1990_10_Problem_Figure_1.png}
                    \end{figure}
                    
$\text{(A)}\ \text{P} \qquad \text{(B)}\ \text{Q} \qquad \text{(C)}\ \text{R} \qquad \text{(D)}\ \text{S} \qquad \text{(E)}\ \text{T}$


	\item (\textit{AMC 8 1990}) The numbers on the faces of this cube are consecutive whole numbers.  The sums of the two numbers on each of the three pairs of opposite faces are equal.  The sum of the six numbers on this cube is


 
                    \begin{figure}[H]
                        \centering
                        \includegraphics[width=0.5\textwidth]{Problem/AMC_8_AJHSME_1990_11_Problem_Figure_1.png}
                    \end{figure}
                    
$\text{(A)}\ 75 \qquad \text{(B)}\ 76 \qquad \text{(C)}\ 78 \qquad \text{(D)}\ 80 \qquad \text{(E)}\ 81$


	\item (\textit{AMC 8 1990}) There are twenty-four $4$-digit numbers that use each of the four digits $2$, $4$, $5$, and $7$ exactly once.  Listed in numerical order from smallest to largest, the number in the $17\text{th}$ position in the list is


 $\text{(A)}\ 4527 \qquad \text{(B)}\ 5724 \qquad \text{(C)}\ 5742 \qquad \text{(D)}\ 7245 \qquad \text{(E)}\ 7524$


	\item (\textit{AMC 8 1990}) One proposal for new postage rates for a letter was $30$ cents for the first ounce and $22$ cents for each additional ounce (or fraction of an ounce).  The postage for a letter weighing $4.5$ ounces was


 $\text{(A)}\ \text{96 cents} \qquad \text{(B)}\ \text{1.07 dollars} \qquad \text{(C)}\ \text{1.18 dollars} \qquad \text{(D)}\ \text{1.20 dollars} \qquad \text{(E)}\ \text{1.40 dollars}$


	\item (\textit{AMC 8 1990}) A bag contains only blue balls and green balls.  There are $6$ blue balls.  If the probability of drawing a blue ball at random from this bag is $\frac{1}{4}$, then the number of green balls in the bag is


 $\text{(A)}\ 12 \qquad \text{(B)}\ 18 \qquad \text{(C)}\ 24 \qquad \text{(D)}\ 30 \qquad \text{(E)}\ 36$


	\item (\textit{AMC 8 1990}) The area of this figure is $100\text{ cm}^2$.  Its perimeter is 


 
                    \begin{figure}[H]
                        \centering
                        \includegraphics[width=0.5\textwidth]{Problem/AMC_8_AJHSME_1990_15_Problem_Figure_1.png}
                    \end{figure}
                    


 [figure consists of four identical squares]

 $\text{(A)}\ \text{20 cm} \qquad \text{(B)}\ \text{25 cm} \qquad \text{(C)}\ \text{30 cm} \qquad \text{(D)}\ \text{40 cm} \qquad \text{(E)}\ \text{50 cm}$


	\item (\textit{AMC 8 1990}) $1990-1980+1970-1960+\cdots -20+10 =$
$\text{(A)}\ -990 \qquad \text{(B)}\ -10 \qquad \text{(C)}\ 990 \qquad \text{(D)}\ 1000 \qquad \text{(E)}\ 1990$


	\item (\textit{AMC 8 1990}) A straight concrete sidewalk is to be $3$ feet wide, $60$ feet long, and $3$ \textbf{inches} thick.  How many cubic yards of concrete must a contractor order for the sidewalk if concrete must be ordered in a whole number of cubic yards?


 $\text{(A)}\ 2 \qquad \text{(B)}\ 5 \qquad \text{(C)}\ 12 \qquad \text{(D)}\ 20 \qquad \text{(E)}\ \text{more than 20}$


	\item (\textit{AMC 8 1990}) Each corner of a rectangular prism is cut off.  Two (of the eight) cuts are shown.  How many edges does the new figure have?


 
                    \begin{figure}[H]
                        \centering
                        \includegraphics[width=0.5\textwidth]{Problem/AMC_8_AJHSME_1990_18_Problem_Figure_1.png}
                    \end{figure}
                    
$\text{(A)}\ 24 \qquad \text{(B)}\ 30 \qquad \text{(C)}\ 36 \qquad \text{(D)}\ 42 \qquad \text{(E)}\ 48$


	\item (\textit{AMC 8 1990}) There are $120$ seats in a row.  What is the fewest number of seats that must be occupied so the next person to be seated must sit next to someone?


 $\text{(A)}\ 30 \qquad \text{(B)}\ 40 \qquad \text{(C)}\ 41 \qquad \text{(D)}\ 60 \qquad \text{(E)}\ 119$


	\item (\textit{AMC 8 1990}) The annual incomes of $1,000$ families range from $8200$ dollars to $98,000$ dollars.  In error, the largest income was entered on the computer as $980,000$ dollars.  The difference between the mean of the incorrect data and the mean of the actual data is


 $\text{(A)}\ \text{882 dollars} \qquad \text{(B)}\ \text{980 dollars} \qquad \text{(C)}\ \text{1078 dollars} \qquad \text{(D)}\ \text{482,000 dollars} \qquad \text{(E)}\ \text{882,000 dollars}$


	\item (\textit{AMC 8 1990}) A list of $8$ numbers is formed by beginning with two given numbers.  Each new number in the list is the product of the two previous numbers.  Find the first number if the last three are shown:

 \[\text{\underline{\hspace{3 mm}?\hspace{3 mm}}\hspace{1 mm},\hspace{1 mm} \underline{\hspace{7 mm}}\hspace{1 mm},\hspace{1 mm} \underline{\hspace{7 mm}}\hspace{1 mm},\hspace{1 mm} \underline{\hspace{7 mm}}\hspace{1 mm},\hspace{1 mm} \underline{\hspace{7 mm}}\hspace{1 mm},\hspace{1 mm}\underline{\hspace{2 mm}16\hspace{2 mm}}\hspace{1 mm},\hspace{1 mm}\underline{\hspace{2 mm}64\hspace{2 mm}}\hspace{1 mm},\hspace{1 mm}\underline{\hspace{1 mm}1024\hspace{1 mm}}}\]
$\text{(A)}\ \frac{1}{64} \qquad \text{(B)}\ \frac{1}{4} \qquad \text{(C)}\ 1 \qquad \text{(D)}\ 2 \qquad \text{(E)}\ 4$


	\item (\textit{AMC 8 1990}) Several students are seated at a large circular table.  They pass around a bag containing $100$ pieces of candy.  Each person receives the bag, takes one piece of candy and then passes the bag to the next person.  If Chris takes the first and last piece of candy, then the number of students at the table could be


 $\text{(A)}\ 10 \qquad \text{(B)}\ 11 \qquad \text{(C)}\ 19 \qquad \text{(D)}\ 20 \qquad \text{(E)}\ 25$


	\item (\textit{AMC 8 1990}) The graph relates the distance traveled [in miles] to the time elapsed [in hours] on a trip taken by an experimental airplane.  During which hour was the average speed of this airplane the largest?


 
                    \begin{figure}[H]
                        \centering
                        \includegraphics[width=0.5\textwidth]{Problem/AMC_8_AJHSME_1990_23_Problem_Figure_1.png}
                    \end{figure}
                    
$\text{(A)}\ \text{first (0-1)} \qquad \text{(B)}\ \text{second (1-2)} \qquad \text{(C)}\ \text{third (2-3)} \qquad \text{(D)}\ \text{ninth (8-9)} \qquad \text{(E)}\ \text{last (11-12)}$


	\item (\textit{AMC 8 1990}) Three $\Delta$'s and a $\diamondsuit$ will balance nine $\bullet$'s.  One $\Delta$ will balance a $\diamondsuit$ and a $\bullet$.


 
                    \begin{figure}[H]
                        \centering
                        \includegraphics[width=0.5\textwidth]{Problem/AMC_8_AJHSME_1990_24_Problem_Figure_1.png}
                    \end{figure}
                    


 How many $\bullet$'s will balance the two $\diamondsuit$'s in this balance?


 
                    \begin{figure}[H]
                        \centering
                        \includegraphics[width=0.5\textwidth]{Problem/AMC_8_AJHSME_1990_24_Problem_Figure_2.png}
                    \end{figure}
                    
$\text{(A)}\ 1 \qquad \text{(B)}\ 2 \qquad \text{(C)}\ 3 \qquad \text{(D)}\ 4 \qquad \text{(E)}\ 5$


	\item (\textit{AMC 8 1990}) How many different patterns can be made by shading exactly two of the nine squares?  Patterns that can be matched by flips and/or turns are not considered different.  For example, the patterns shown below are \textbf{not} considered different.


 
                    \begin{figure}[H]
                        \centering
                        \includegraphics[width=0.5\textwidth]{Problem/AMC_8_AJHSME_1990_25_Problem_Figure_1.png}
                    \end{figure}
                    
$\text{(A)}\ 3 \qquad \text{(B)}\ 6 \qquad \text{(C)}\ 8 \qquad \text{(D)}\ 12 \qquad \text{(E)}\ 18$


\end{enumerate}